\section{Scope and further questions}\label{sec:discussion}

The results address three different decisions. A minimax formula gives the
worst rounding error before the relaxed input is known. An exact instance
algorithm finds the best schedule once the input is supplied. A coarsening
certificate bounds the loss caused by reducing the allowed switching grid.
Together they separate switch-budget error, grid restriction, and input
representation error, and the minimax analysis imposes no smoothness on the
measurable control.

Two separate Lean developments, listed in the supplement's README, verify
selected results. The \texttt{SwitchingControl} package proves the exact
five-, six-, and seven-cell three-mode values and the continuous value
$F_{3,2}(T)=T/5$, including measurable inputs; it does not cover the chamber
counts, the optimal-switch histograms, or the converse direction of
Proposition~\ref{prop:floor-automaton}. The \texttt{GridSwitching} package
proves the arbitrary-grid one-switch formula, the transfer and coarsening
certificates, and the supporting unrestricted instance algorithms. Their
coverage and verification records specify the proved statements and retained
checks. Neither verifies the whole manuscript, the Python implementations,
timing or bit-cost claims, operation counts beyond candidate-set sizes, or
bibliographic novelty.

Three mathematical extensions remain substantial. First, the exact
five-block one-sided reach bound for general mode count is open. The
exclusion identity in Appendix~\ref{sec:higher-reach} identifies the missing
weighted three-block inequality. The negative event assignment there is not a
control trajectory, and chronological monotonicity removes that particular
obstruction, but a certificate for one chamber covers neither the remaining
orderings nor the full reach bound. The general and seeded results of
Section~\ref{sec:general-budgets} do not depend on this question. A solution
would improve the higher-budget minimax formulas and the second-order
mode-count gap.

Second, the regime with at least as many activation blocks as modes has a
different structure. The three-mode one-sided obstruction rules out a direct
extension of the distinct-mode method, yet the exact value $F_{3,2}(T)=T/5$
and the range $T/7\le F_{3,3}(T)\le T/6$ show that finite-grid chamber methods
can still improve the boundary theory. The floor-history construction is
valid for every grid length, but the enumerations here stop at seven cells.
The six exceptional seven-cell chambers have explicit repairs; their existence
also disproves the tempting unit-error extension on $2s+1$ cells. Determining
$F_{3,3}$ and a general three-mode floor-history switch law requires further
arguments. For other small-mode regimes, the classical
bound~\eqref{eq:classical-small-mode} remains a useful comparison.

Third, constraints beyond a switch budget change the approximation question.
The exact subset algorithm handles its stated per-mode minimum dwell rules,
but a transition graph couples successive labels and is outside that unary
assignment reduction. The odd-grid dwell example shows that unrestricted
coarsening certificates can fail even under arbitrarily fine refinement.
Useful constrained certificates would need assumptions that ensure feasible
nearby switching times, or a constraint-aware refinement construction.

The algorithms already support exact small-budget rounding and certified
additive approximation with rational data. A broader computational study
could compare them with established graph and mixed-integer methods at
matched budgets and objectives, and then evaluate the resulting trajectories
under specified dynamics. Such a study would address application performance,
not the scope of the discrepancy theorems.
